% !PRINTED none
\begin{center}
\gesperrt{Fünfundzwanzigster Abschnitt.}

Funktionale.

\rule{4cm}{0.4pt}
\end{center}

\begin{center}
§ 177. \gesperrt{Rationale Funktionale.}
\end{center}

Wir übertragen nun den Begriff des Funktionals, der uns im
17. Abschnitt des zweiten Bandes zur Begründung der Theorie
der algebraischen Zahlen gedient hat, auf die algebraischen Funktionen.

Wir adjungieren also dem Körper $\Omega$ beliebige Variable und
rechnen damit nach den Regeln der Buchstabenrechnung. Es
entsteht so ein erweiterter Körper $\overline{\Omega}$, dessen Elemente \gesperrt{Funktionale}
 heißen. Der Körper $\overline{\Omega}$ selbst heißt der \gesperrt{Funktionalkörper}.
 Die Variablen sind hier nicht im Sinne der Analysis
als Zeichen für veränderliche Zahlen, sondern als bloße Rechnungssymbole
 aufzufassen, und sind wohl zu unterscheiden von den
Variablen $z, \theta$ des Körpers $\Omega$. Wir wollen diese Hilfsvariablen
die \gesperrt{Funktionalvariablen} nennen.

Aus dem Körper $Z$ der rationalen Funktionen von $z$ entsteht
durch Adjunktion der Variablen der \gesperrt{Körper $\overline{Z}$ der rationalen
Funktionale}.

Eine ganze rationale Funktion der Variablen:

\[
\Phi(u, v, w, \ldots)
\]

mit ganzen rationalen Funktionen von $z$ als Koeffizienten heißt
eine ganze Funktion in $Z$. Der größte gemeinschaftliche Teiler
der Koeffizienten von $\Phi$ heißt \gesperrt{der Teiler der Funktion}, und
die Funktion heißt \gesperrt{primitiv}, wenn die Koeffizienten keinen gemeinschaftlichen
 Teiler haben. Die primitiven rationalen Funktionen
 und ihre Quotienten werden auch \gesperrt{Einheiten im Körper}
$\overline{Z}$ genannt.
